% Independent D1 coaxial spectral reference. 2026-11-06; D1 remains open.
\section{D1 独立共轴分层 Helmholtz 物理直接参考：谱域传播与端口--场一致性}

\paragraph{方法迁移，而不是论文数字拟合。}
本节谱域 Helmholtz 分离、衰减分支、源跳跃、层界面 DtN、
复端口与 Joule 功率是为了学习并提取成熟计算方法。
原论文表格对照只作为受限轴对称模型的历史附加证据，
\textbf{不是}任意 SE(3) BFZI 求解的必要验收标准，
也不因个别表格差异而继续为其调参。
真实 BFZI 任意位姿和两个绝缘盒需要独立三维
Maxwell--MQS $A$--$\phi$ 保结构参考及
TEAM--7 证据，详见
\texttt{docs/research/bfzi\_uwpt\_electrothermal/METHOD\_TRANSFER\_D1\_TO\_D2.md}。

\label{sec:bfzi-d1-layered-helmholtz}

\paragraph{角色及严格适用范围。}
本节是物理直接法锁定规范（第~\ref{sec:bfzi-physical-direct-locked}~节）的
D1 半解析标准问题，\emph{不是} BFZI 任意位姿和绝缘双盒的解析解。
参照 Shafiei 等~\cite{shafiei2025uwpt}，两组圆形共轴薄环位于平面
$z=0$ 和 $z=d$，各向同性常数磁导率取 $\mu_0$，三个
沿 $z$ 的平面层横向无限，材料参数为
$(\sigma_j,\varepsilon_j)$，$j=1,2,3$。
所有薄环同向端口，峰值相量约定 $\exp(+\mathrm i\omega t)$；
相对于真实 BFZI 连续闭合馈回流，此处只用封闭圆环的标准求解基准。

\subsection{复传播常数及界面 Dirichlet--Neumann 消元}
取 Hankel 波数 $k\geq0$、$J_1$ 为一阶 Bessel 函数，
由均匀材料层中的 $\mathrm A_\varphi$ 或等价的
$\mathrm E_\varphi=-\mathrm i\omega \mathrm A_\varphi$ 满足的
Helmholtz 方程，可取
\begin{equation}
u_j(k)^2=k^2+\mathrm i\omega\mu_0\sigma_j
                 -\omega^2\mu_0\varepsilon_j,\qquad
q_j=\frac{u_j}{\mu_0},\qquad t=e^{-u_2d}.
\end{equation}
平方根取保证半空间衰减/出射的分支。
D1 对论文完整复现保留位移电流项；对 MQS 验证显式去掉
$-\omega^2\mu_0\varepsilon_j$，二者不可无条件混称。

若在两个源面定义谱域环向位势的迹
$(F_0,F_d)^T$，由两个源平面的跳跃条件、
上下半空间出射条件和中间层双向传播，得到
\begin{equation}
\begin{bmatrix}F_0\\F_d\end{bmatrix}
=\bm G_{\rm int}(k)
\begin{bmatrix}S_1(k)I_1\\S_2(k)I_2\end{bmatrix},
\quad
S_p(k)=\sum_{n=1}^{N_p}a_{pn}J_1(ka_{pn}),
\end{equation}
其中
\begin{align}
D &= (q_2+q_1)(q_2+q_3)-(q_2-q_1)(q_2-q_3)t^2,\\
G_{00} &=\frac{(q_2+q_3)+(q_2-q_3)t^2}{D},&
G_{0d}&=\frac{2q_2t}{D},&
G_{dd}&=\frac{(q_2+q_1)+(q_2-q_1)t^2}{D},\\
\bm G_{\rm int}&=
\begin{bmatrix}G_{00}&G_{0d}\\G_{0d}&G_{dd}\end{bmatrix}.
\end{align}
这是一种不含增长指数 $e^{+u_2d}$ 的等价数值表示。
上下介质完全相同时应严格退化为
$G_{00}=G_{dd}=\mu_0/(2u)$ 和
$G_{0d}=\mu_0 e^{-ud}/(2u)$。
源面法向导数跳跃在谱域满足
\begin{equation}
\frac{\partial_zF(0^-)-\partial_zF(0^+)}{\mu_0}=S_1I_1,\quad
\frac{\partial_zF(d^-)-\partial_zF(d^+)}{\mu_0}=S_2I_2,
\end{equation}
并对两个独立源列分别验证。这里的 $F$ 为谱域振幅；
实际场按 $A_\varphi(r,z)=\int_0^\infty
 kJ_1(kr)F(k,z)\,\mathrm dk$ 重建。

\subsection{复增量端口与空间物理量}
对零半径线源，真空自感发散，因此只在\emph{相同薄环集合}上
定义有限的介质反作用增量
\begin{equation}
\boxed{
\Delta Z_{pq}
= 2\pi\mathrm i\omega
\int_0^\infty k\, S_p(k)\,
\left[G_{pq}(k)-G^{\rm vac}_{pq}(k)\right]\,
S_q(k)\,\mathrm dk.}
\end{equation}
这不会产生有限铜截面的绝对自阻抗，也不含补偿电容、
整流器或负载。复感定义
$\Delta L_{pq}=\Delta Z_{pq}/(\mathrm i\omega)$，
正的海水耗散对应其负虚部，
与 Wang 等~\cite{wang2025marinewpt} 所用的
场--路复参数思想相容，但 Wang 具体串--并联外电路
须独立按 KCL/KVL 组合，不能简单当作串联对角阻抗。

通过谱域重建（仅在非源面有效）：
\begin{equation}
E_\varphi=-\mathrm i\omega A_\varphi,\qquad
B_r=-\partial_z A_\varphi,\qquad
B_z=\frac1r\partial_r(rA_\varphi),\qquad
J_\varphi=\sigma_j E_\varphi,\qquad
q_j=\tfrac12\sigma_j|E_\varphi|^2.
\end{equation}
在这个共轴环向纯模态问题中，$\varphi$ 的附加梯度
不产生环向分量；一般三维 BFZI \emph{不许}据此省略
$A$--$\phi$ 中的电标势及规范条件。

\paragraph{总场与海水反作用场不可混淆。}
令 $E^{\rm sec}_\varphi=E^{\rm med}_\varphi-
E^{\rm vac}_\varphi$，两项严格在同一线源、频率和
端口相量下计算。Shafiei 等所列表~2 的六组电场峰值
在数值上与本模型的 $|E^{\rm sec}_\varphi|$ 峰值高度吻合，
而不与源附近的 $|E^{\rm med}_\varphi|$ 相符；
但作者未显式给出这一等式作为表项定义，
因此这只是目前最有数据支持的\emph{复现解释}，
不是独立精度认证。
海水热源必须计算
$q_j=\tfrac12\sigma_j|E^{\rm med}_\varphi|^2$，
\textbf{禁止}以 $\sigma_j|E^{\rm sec}_\varphi|^2/2$
冒充真实功率。

在 MQS 下，Bessel--Parseval 恒等式给出可执行的
端口--体积功率等价检查：
\begin{equation}
P_{\rm sea}
=\frac12\int_{\Omega_{\rm sea}}\sigma|\bm E^{\rm med}|^2\,\mathrm dV
=\pi\omega^2\int_0^\infty k\sum_{j=1}^3
\sigma_j\int_{\text{layer }j}|F(k,z)|^2\,\mathrm dz\,\mathrm dk
=\tfrac12\bm I^{\rm H}\Herm{\Delta\bm Z}\bm I.
\end{equation}
两侧半空间的 $z$ 积分可解析求出，有限厚度层单独数值积分；
这仍是同一半解析模型的内部场--端口一致性检查。
保留位移电流时必须分账远场辐射，
不能把真空扣除后增量阻抗的全部实部自动定义成海水 Joule 损耗。

\subsection{实算结果和 D1 状态}
对原论文表~1 的 $N_1/N_2=15/17$、内径
$170/165\,$mm、导线之间净距 $1\,$mm、
导线直径 $1.5\,$mm，
按圆环中心径向间距 $2.5\,$mm 进行 D1 薄环标准件计算。
无限与半无限海水在 $85/400\,$kHz 的附加
自阻和互阻复现原文表~2 的舍入精度；
有限海水层在 $400\,$kHz 下的三项电阻仍低于论文
解析栏约 $4$--$5\%$，且对谱域截断较稳定。
三种介质在两频率的海水反作用电场离散峰值与原文表~2
相差约 $0.33\%$--$4.15\%$；该比较的观测量
定义尚待原论文/作者确认，\textbf{不作事后设门槛的 D1 PASS}。
全部复现参数、空间网格、数据和可运行命令见
\texttt{docs/research/bfzi\_uwpt\_electrothermal/D1\_SHAFIEI\_SPECTRAL\_REPRODUCTION.md}。
D1 仍 OPEN，D2 独立三维 $A$--$\phi$/TEAM--7 尚未验收；
旧 G1 五模态 FAIL、正式神经教师冻结。
